\section{Convention 2: Which Hits Count}
\label{sec:matching}
The conventions an evaluation fixes are not ours alone. PrefixSim, the open simulator of a recent study of
fourteen eviction policies on production prefix-cache traces~\cite{liu2026fancyeviction}, counts a cached
block as a hit even when an earlier block of the same request missed; its documentation says so, and notes
that vLLM and SGLang stop matching at a request's first miss. Under that count, a compute-aware policy that
evicts cheap, shallow, long-idle blocks (RandomCompute) saves 10.4 points more prefill compute than LRU at the
study's tightest cache, its headline. We added a second count to the same match, hits up to each request's
first miss, leaving the cache's behaviour unchanged, and checked it against the simulator's own record of each
request's missing runs: the two agree to $10^{-6}$ for every policy recorded. Before running, we registered
that the headline survives if the lead stays at 5 points or more and does not if it falls below one
(\Cref{tab:matching}).

\begin{table}[t]
\centering\footnotesize\setlength{\tabcolsep}{3.5pt}
\begin{tabular}{@{}lrrrr@{}}
\toprule
& \multicolumn{2}{c}{freeinference} & \multicolumn{2}{c}{chutes} \\
\cmidrule(lr){2-3}\cmidrule(l){4-5}
16k blocks & all hits & to first miss & all hits & to first miss \\
\midrule
LRU & 53.3 & 53.3 & 60.0 & 60.0 \\
Belady (hits) & 60.9 & 60.9 & 67.6 & 67.6 \\
RandomCompute & 63.8 & \textbf{33.7} & 60.8 & \textbf{45.7} \\
BeladyCompute (offline) & 73.5 & \textbf{36.1} & 68.7 & 62.8 \\
\bottomrule
\end{tabular}
\caption{Percent of prefill compute saved by PrefixSim at its published commit on its two public traces,
16k blocks of Qwen3-Coder-30B's KV (24\,GiB), its cost model, counting every resident block as a hit or
only the hits before a request's first miss.}
\label{tab:matching}
\end{table}

\noindent Counted the way the engines match, RandomCompute's lead becomes a loss of 19.6 points on
freeinference and 14.3 on chutes, and even the compute-aware offline oracle falls below LRU on freeinference.
LRU and the hit-count Belady lose nothing: a request that reads a block reads its whole prefix, so recency
never keeps a block whose parent it has evicted. The compute-aware policies evict the cheap blocks at the front
of a prefix and strand the expensive ones behind them. The partial-node variant the study proposes cuts holes
per request from 9.3 to 1.06, close to LRU's 0.93, yet loses as much (20.9\% of blocks are hits only behind a
hole, against 21.6\%): one hole early in a long prefix is enough. The study is explicit about holes and
proposes an engine change that recomputes them in place; its headline holds for such an engine and reverses on
today's. The difference shrinks as the cache grows, to 1.5 points at 256k blocks. Which hits count is one more
convention an evaluation fixes, and here it decides the sign of the result.
